\section{Aplicación de SSDLC a Workcodile-dev}

Un \textbf{Secure Software Development Life Cycle (SSDLC)} es un SDLC tradicional con actividades de seguridad integradas en cada fase. La seguridad no se revisa al final, sino que es un criterio transversal que acompaña todo el proceso de construcción del software \cite{nist800218}.

\subsection{Fase 1: Requisitos de Seguridad}

Aquí se responde la pregunta: \textit{¿Qué debemos proteger?}

\subsubsection{Activos de información de Workcodile-dev}

\begin{table}[h]
\centering
\small
\resizebox{\columnwidth}{!}{%
\begin{tabular}{|l|l|l|}
\hline
\textbf{Activo} & \textbf{Tipo} & \textbf{Ubicación} \\
\hline
Base de datos MongoDB & Datos & Backend (contenedor) \\
Credenciales de usuario & Sensibles & MongoDB (bcrypt hash) \\
Sesiones JWT & Tokens & Memoria del servidor \\
Archivos subidos & Datos & MinIO (contenedor) \\
Variables .env & Secretos & Servidor host \\
Código fuente & Propiedad & Repositorio Git \\
\hline
\end{tabular}%
}
\caption{Activos de información del proyecto}
\label{tab:activos}
\end{table}

\subsubsection{Clasificación de la información}

\begin{itemize}
    \item \textbf{Pública:} Cursos disponibles, perfil público de usuarios.
    \item \textbf{Interna:} Código fuente, configuraciones de Docker, estructura de la base de datos.
    \item \textbf{Confidencial:} Correos electrónicos de usuarios, contraseñas hasheadas, tokens JWT.
    \item \textbf{Restringida:} Variables de entorno con contraseñas de servicios (Gmail, MongoDB, MinIO).
\end{itemize}

\subsubsection{Requisitos de seguridad (CIA + T)}

\begin{itemize}
    \item \textbf{Confidencialidad:} Solo usuarios autenticados acceden a sus datos.
    \item \textbf{Integridad:} Los datos de publicaciones no deben ser modificados sin autorización.
    \item \textbf{Disponibilidad:} La aplicación debe estar operativa para los usuarios.
    \item \textbf{Trazabilidad:} Las acciones relevantes deben poder registrarse y auditarse.
\end{itemize}

\subsection{Fase 2: Diseño Seguro}

Aquí se responde la pregunta: \textit{¿Cómo podrían atacarnos?}

\subsubsection{Threat Modeling con STRIDE}

Aplicando la metodología STRIDE a la arquitectura de Workcodile-dev:

\begin{table}[h]
\centering
\small
\resizebox{\columnwidth}{!}{%
\begin{tabular}{|l|p{4.5cm}|p{4.5cm}|}
\hline
\textbf{Amenaza} & \textbf{Ejemplo} & \textbf{Control} \\
\hline
Spoofing & Suplantación con JWT falso & Validación firma JWT \\
Tampering & Modificación datos en tránsito & HTTPS, validación ObjectId \\
Repudiation & Usuario niega publicación & Logging con timestamp \\
Info Disclosure & Exposición de .env & .gitignore excluye .env \\
DoS & Inundación a la API & Rate limiting \\
Elevation of Priv. & Acceso a rutas admin & Middleware por roles \\
\hline
\end{tabular}%
}
\caption{Análisis STRIDE de Workcodile-dev}
\label{tab:stride}
\end{table}

\subsubsection{Diagrama de Flujo de Datos (DFD)}

La arquitectura de Workcodile-dev presenta los siguientes flujos principales de datos:

\begin{enumerate}
    \item \textbf{Flujo de autenticación:} Usuario → Frontend (React) → Backend (Express) → MongoDB → JWT token.
    \item \textbf{Flujo de publicaciones:} Usuario → Frontend → Backend → MinIO (archivo) + MongoDB (metadatos).
    \item \textbf{Flujo de correo:} Backend → Nodemailer → Gmail SMTP.
\end{enumerate}

Cada frontera entre componentes representa un límite de confianza que requiere controles adicionales de validación y sanitización.

\subsection{Fase 3: Desarrollo Seguro}

En esta fase se establecen las reglas de codificación segura para el proyecto. El control A.14.2 de ISO 27001 exige que la organización defina procesos de codificación segura que incluyan \cite{iso27001}:

\begin{itemize}
    \item Principios de codificación segura aprobados.
    \item Requisitos de seguridad para desarrollo interno y externo.
    \item Procedimientos para prevenir vulnerabilidades conocidas.
    \item Uso de herramientas de desarrollo configuradas de forma segura.
    \item Entornos de desarrollo controlados y protegidos.
\end{itemize}

Estas prácticas se detallan en la Sección~\ref{sec:secure-coding} del presente informe.

\subsection{Fase 4: Pruebas de Seguridad}

El control A.14.3 de ISO 27001 exige que las pruebas de seguridad formen parte del proceso normal de pruebas del sistema \cite{iso27001}. Las verificaciones deben incluir:

\begin{itemize}
    \item Funciones de autenticación y autorización.
    \item Controles de acceso.
    \item Gestión de sesiones.
    \item Registro y auditoría de eventos.
    \item Configuraciones de seguridad.
\end{itemize}

Las actividades de verificación incluyen SAST, DAST, revisiones de código y escaneo de vulnerabilidades. Esto se detalla en la Sección~\ref{sec:security-testing}.

\subsection{Fase 5: Despliegue Seguro}

Para Workcodile-dev, el despliegue se realiza mediante Docker Compose. Los controles de seguridad incluyen:

\subsubsection{Hardening de contenedores}
\begin{itemize}
    \item Imágenes mínimas (node:alpine en lugar de node:full).
    \item Usuarios no root en contenedores de producción.
    \item Configuración restrictiva de puertos.
    \item Eliminación de servicios innecesarios.
\end{itemize}

\subsubsection{Gestión de secretos}
\begin{itemize}
    \item Variables de entorno en archivo \texttt{.env} (excluido de Git).
    \item Archivo \texttt{.env.example} como plantilla versionada.
    \item Nunca secretos hardcodeados en código fuente.
\end{itemize}

\subsubsection{Seguridad CI/CD}
\begin{itemize}
    \item Control de acceso a repositorios.
    \item Protección de ramas principales mediante revisiones.
    \item Validación automática con pruebas antes de cada integración.
    \item Ejecución de análisis de seguridad dentro del pipeline.
    \item Segregación entre entornos de desarrollo, pruebas y producción.
\end{itemize}

\subsection{Fase 6: Mantenimiento}

Las actividades de mantenimiento de seguridad para Workcodile-dev incluyen:

\begin{itemize}
    \item Actualización periódica de dependencias (\texttt{npm audit}).
    \item Corrección de vulnerabilidades identificadas.
    \item Monitoreo continuo de riesgos y amenazas.
    \item Revisión de configuraciones de seguridad.
    \item Revisión de permisos y controles de acceso.
    \item Evaluación de nuevas vulnerabilidades publicadas.
\end{itemize}
